% !PRINTED 423
\weberpage{423}
%<<CONT
hat es einen ungraden quadratischen Primteiler $p^2$, der nicht in $y^2$
aufgeht, und es ist
\[
(y+4\lambda p)^2 + m x^2 = 4M^2 + 8\lambda y p + 16\lambda^2 p^2,
\]
und diese Zahl ist, wenn $\lambda$ nicht durch $p$ teilbar ist, zwar durch
$p$, aber nicht durch $p^2$ teilbar, und kann also kein Quadrat sein.
Überdies kann man über $\lambda$ noch so verfügen, daß $y+4\lambda p$ mit
$x$ keinen ungeraden gemeinsamen Teiler hat, daß also die Bedingungen
§ 114 erfüllt sind. Demnach genügt $j(\omega)$ für jedes $m$ einer
Gleichung $F_n(u,u)$, in der $C=\pm 1$ ist, und daraus folgt nach Bd. II,
§ 149:

VI. \gesperrt{Die Klassenvarianten sind ganze algebraische Zahlen.} % [sic] Der Druck hat „Klassenvarianten“; gemeint sind die Klasseninvarianten.

\begin{center}
\textbf{\S\,116. \gesperrt{Klassenzahlrelationen.}}
\end{center}

Wenn man den Grad, wie er sich hiernach für beide Teile der Gleichung
(10), § 115 ergibt, gleich setzt, so erhält man die Formel
\weq{1}{N = hk + h'k' + h''k'' + \cdots,}
\weq{2}{\sum h k = 2\sum^{\partial} \frac{\partial}{e}\varphi(e) + \varphi(\sqrt{n}),}
worin $\varphi(\sqrt{n})=0$ zu setzen ist, wenn $n$ kein Quadrat ist, und
$N$ durch (13), (14) oder (15) des vorigen Paragraphen bestimmt ist.
Dies ist die \gesperrt{Kroneckersche Klassenzahlrelation}, der wir noch eine etwas
bequemere Form geben wollen\footnote{Die Klassenzahlrelationen, von
denen die hier abgeleitete nur der einfachste Fall ist, sind von
\gesperrt{Kronecker} entdeckt (Crelles Journal, Bd. 57) und von
\gesperrt{Gierster} (Mathematische Annalen, Bd. 21, 22) und
\gesperrt{Hurwitz} (ebenda Bd. 25) bedeutend verallgemeinert.}.

Wir fassen neben $n$ alle Werte $n'$ ins Auge, die aus $n$ durch
Fortheben eines quadratischen Faktors entstehen, bilden für sie die
Gleichung (2) und addieren alle so gewonnenen Resultate. Dabei ist nur,
falls $n$ ein Quadrat ist, der Wert $n'=1$ auszuschließen, weil
$F_1(u,u)=0$ ist.

Es sei also $\delta^2$ irgend ein von $n$ selbst verschiedener
quadratischer Faktor von $n$ und
\[
n = n'\delta^2.
\]
Zerlegen wir $n'$ in zwei Faktoren $n' = a'\partial'$ und bezeichnen mit
$e'$ den größten gemeinschaftlichen Teiler von $\partial'$ und $a'$, so
ist
%CONT>>
